\documentclass[a4paper,10pt]{article}
\usepackage{polski}
\usepackage[latin2]{inputenc}
\usepackage{amsmath}
\usepackage{amsfonts}
\usepackage{amssymb}
\usepackage{slashed}
\usepackage{mathrsfs} 
\usepackage{bm}

\usepackage{geometry}
\geometry{verbose,a4paper,tmargin=2.5cm,bmargin=2.5cm,lmargin=2.5cm,rmargin=2.5cm}

%opening
\title{Integration scheme for GPE}
\author{Gabriel Wlaz�owski}

\begin{document}

\maketitle

% \begin{abstract}
% 
% \end{abstract}

\section{GPE equation}

\begin{equation}
(i\alpha - \beta)\frac{\partial\Psi(\vec{r}, t)}{\partial t} =  \left(  -\frac{1}{4}\nabla^2 + 2\frac{\delta\mathcal{E}(n,t)}{\delta n} + 2 V_{ext}(\vec{r}, t)\right) \Psi(\vec{r}, t)
\end{equation} 
where $\Psi$ is collective wave function of system, $\mathcal{E}$ is energy density functional and $V_{ext}$ is external potential.\\\\
Note:
\begin{enumerate}
 \item $\alpha=1$ and $\beta=0$ corresponds to real time evolution.
 \item $\alpha=0$ and $\beta=1$ corresponds to imaginary time evolution. 
\end{enumerate}

For unitary Fermi gas ($m=\hbar=1$):

\begin{equation}
\mathcal{E}(n) = 0.37 \dfrac{3}{10}n(3\pi^2 n)^{2/3}
\end{equation} 
and
\begin{equation}
\frac{\delta\mathcal{E}(n)}{\delta n} = 0.37 \dfrac{1}{2} (3\pi^2 n)^{2/3}
\end{equation}
where $n(\vec{r}, t)=2|\Psi(\vec{r}, t)|^{2}$.

Energy is computed as:
\begin{equation}
 E = E_{kin} + E_{int} + E_{ext},
\end{equation} 
where:
\begin{equation}
 E_{kin} = \int \Psi^{*}(\vec{r})\dfrac{1}{4}\nabla^2\Psi(\vec{r})\,d^{3}\vec{r}
\end{equation} 
\begin{equation}
E_{int}=\int\mathcal{E}(n(\vec{r}))\,d^{3}\vec{r}
\end{equation} 
\begin{equation}
E_{ext}=\int V_{ext}(\vec{r})n(\vec{r})\,d^{3}\vec{r}
\end{equation}

\section{ABM integration algorithm}

Target: solve equation:
\begin{equation}
 \dfrac{dy(t)}{dt} = f(y,t).
\end{equation} 

ABM algorithm:
\begin{eqnarray}
 y_k^{(p)} &=& y_{k-1} + \dfrac{23}{12}\Delta t f_{k-1} - \dfrac{16}{12}\Delta t f_{k-2} + \dfrac{5}{12}\Delta t f_{k-3},\\
 y_k &=& y_{k-1} + \dfrac{9}{24}\Delta t f(y_k^{(p)},k\Delta t) +\dfrac{19}{24}\Delta t f_{k-1}-\dfrac{5}{24}\Delta t f_{k-2}+\dfrac{1}{24}\Delta t f_{k-3}, \label{eqn:ABM_corr}
\end{eqnarray} 
where $f_{k}=f(y_k, k\Delta t)$.\\\\
NOTE: Once can improved integration by adding extra ``iteration'' of Eq.(\ref{eqn:ABM_corr}), i.e. we execute it again with $y_k$ instead of $y_k^{(p)}$.\\\\
It is possible to use only 4 wave function buffers called: $y_{k-1}$, $f_{k-1}$, $f_{k-2}$ and $f_{k-3}$.
\\\\
ALGORITHM: (operations listed within single step have to be executed simultaneously)\\
\begin{enumerate}
 \item 
 \begin{equation}
  \left\lbrace 
  \begin{array}{lll}
   y_{k-1} & \leftarrow & y_{k-1} + \dfrac{23}{12}\Delta t f_{k-1} - \dfrac{16}{12}\Delta t f_{k-2} + \dfrac{5}{12}\Delta t f_{k-3},\\
   f_{k-3} & \leftarrow & y_{k-1} +\dfrac{19}{24}\Delta t f_{k-1}-\dfrac{5}{24}\Delta t f_{k-2}+\dfrac{1}{24}\Delta t f_{k-3}
  \end{array}
  \right. 
 \end{equation} 
  \item Compute densities from $y_{k-1}$ and formulate potentials (like $U(\vec{r})$ and $\Delta(\vec{r})$) needed to formulate Hamiltonian. This step requires MPI communication. For densities and potentials use additional buffers, note the are small in comparison to wave function buffers.
  \item $y_{k-1} \leftarrow f(y_{k-1},k\Delta t)$, this is first application of Hamiltonian, as Hamiltonian use potentials formulated in step 2. This step can be done is batched mode. This is numerically intensive step.
  \item $y_{k-1} \leftarrow \dfrac{9}{24}\Delta t y_{k-1} + f_{k-3}$, this step can be merged with step 3. Now buffer $y_{k-1}$ holds wave functions for $t=k\Delta t$.
  \item Compute densities from $y_{k-1}$ and formulate potentials (like $U(\vec{r})$ and $\Delta(\vec{r})$) needed to formulate Hamiltonian. This step requires MPI communication. For densities and potentials use additional buffers, note the are small in comparison to wave function buffers.
  \item Preparation of buffers for next step:
 \begin{equation}
  \left\lbrace 
  \begin{array}{lll}
   f_{k-3} & \leftarrow & f_{k-2},\\
   f_{k-2} & \leftarrow & f_{k-1},\\
   f_{k-1} & \leftarrow & f(y_{k-1},k\Delta t).
  \end{array}
  \right. 
 \end{equation}   
 Note, last operation corresponds to second  application of Hamiltonian, as Hamiltonian use potentials formulated in step 5. This step can be done is batched mode. This is numerically intensive step.
  Finally increase $k$ by one.
\end{enumerate}

\section{ABM2 integration algorithm}

ABM algorithm:
\begin{eqnarray}
 y_k^{(p)} &=& y_{k-1} + \dfrac{55}{24}\Delta t f_{k-1} - \dfrac{59}{24}\Delta t f_{k-2} + \dfrac{37}{24}\Delta t f_{k-3} - \dfrac{9}{24}\Delta t f_{k-4},\\
 y_k &=& y_{k-1} + \dfrac{251}{720}\Delta t f(y_k^{(p)},k\Delta t) +\dfrac{646}{720}\Delta t f_{k-1}-\dfrac{264}{720}\Delta t f_{k-2}+\dfrac{106}{720}\Delta t f_{k-3} - \dfrac{19}{720}\Delta t f_{k-4}, \label{eqn:ABM2_corr}
\end{eqnarray} 
where $f_{k}=f(y_k, k\Delta t)$.\\\\
It is possible to use only 5 wave function buffers called: $y_{k-1}$, $f_{k-1}$, $f_{k-2}$, $f_{k-3}$, $f_{k-4}$.
\\\\
ALGORITHM: (operations listed within single step have to be executed simultaneously)\\
\begin{enumerate}
 \item 
 \begin{equation}
  \left\lbrace 
  \begin{array}{lll}
   y_{k-1} & \leftarrow & y_{k-1} + \dfrac{55}{24}\Delta t f_{k-1} - \dfrac{59}{24}\Delta t f_{k-2} + \dfrac{37}{24}\Delta t f_{k-3} - \dfrac{9}{24}\Delta t f_{k-4},\\
   f_{k-4} & \leftarrow & y_{k-1} +\dfrac{646}{720}\Delta t f_{k-1}-\dfrac{264}{720}\Delta t f_{k-2}+\dfrac{106}{720}\Delta t f_{k-3} - \dfrac{19}{720}\Delta t f_{k-4}
  \end{array}
  \right. 
 \end{equation} 
  \item Compute densities from $y_{k-1}$ and formulate potentials (like $U(\vec{r})$ and $\Delta(\vec{r})$) needed to formulate Hamiltonian. This step requires MPI communication. For densities and potentials use additional buffers, note the are small in comparison to wave function buffers.
  \item $y_{k-1} \leftarrow f(y_{k-1},k\Delta t)$, this is first application of Hamiltonian, as Hamiltonian use potentials formulated in step 2. This step can be done is batched mode. This is numerically intensive step.
  \item $y_{k-1} \leftarrow \dfrac{251}{720}\Delta t y_{k-1} + f_{k-4}$, this step can be merged with step 3. Now buffer $y_{k-1}$ holds wave functions for $t=k\Delta t$.
  \item Compute densities from $y_{k-1}$ and formulate potentials (like $U(\vec{r})$ and $\Delta(\vec{r})$) needed to formulate Hamiltonian. This step requires MPI communication. For densities and potentials use additional buffers, note the are small in comparison to wave function buffers.
  \item Preparation of buffers for next step:
 \begin{equation}
  \left\lbrace 
  \begin{array}{lll}
   f_{k-4} & \leftarrow & f_{k-3},\\
   f_{k-3} & \leftarrow & f_{k-2},\\
   f_{k-2} & \leftarrow & f_{k-1},\\
   f_{k-1} & \leftarrow & f(y_{k-1},k\Delta t).
  \end{array}
  \right. 
 \end{equation}   
 Note, last operation corresponds to second  application of Hamiltonian, as Hamiltonian use potentials formulated in step 5. This step can be done is batched mode. This is numerically intensive step.
  Finally increase $k$ by one.
\end{enumerate}

\end{document}


